\documentclass[border=5pt]{standalone}
% 命名色表必须在 pgfplots 之前声明：pgfplots 内部会先加载 xcolor，
% 之后再 \usepackage[dvipsnames,svgnames]{xcolor} 会触发 Option clash
\PassOptionsToPackage{dvipsnames,svgnames}{xcolor}
\usepackage{pgfplots}
\usepackage{pgf-pie} % 饼图：AI 代码可直接使用 \pie
\pgfplotsset{compat=1.18}
\usepgfplotslibrary{fillbetween}  % 面积图 / 堆叠面积图 / \closedcycle 填充路径
% 注：error bars 是 pgfplots 内置功能（在核心 pgfplots.errorbars.code.tex），不需要单独 \usepgfplotslibrary 加载
\usepackage{amsmath}
\usepackage{amssymb}

% 支持中文
\usepackage{fontspec}
\usepackage{xeCJK}
\setCJKmainfont{SimSun}
\setmainfont{Times New Roman}

\begin{document}

\begin{tikzpicture}\begin{axis}[
    ybar,
    title={各车间产量对比},
    ylabel={产量（台）},
    symbolic x coords={一车间,二车间,三车间},
    xtick=data,
    nodes near coords,
    every node near coord/.append style={font=\scriptsize, fill=none, draw=none, inner sep=1pt, anchor=south},
    legend style={at={(1.03,0.5)}, anchor=west, draw=black, fill=white}
]
\addplot coordinates {(一车间,4520) (二车间,6100) (三车间,3890)};
\legend{产量}
\end{axis}
\end{tikzpicture}

\end{document}